\documentclass[10pt,a4paper]{amsart}
\usepackage[margin=19mm]{geometry}
\usepackage{amsmath,amssymb,mathtools}
\usepackage[scr=rsfs,cal=euler]{mathalfa}
\usepackage{xfrac}
\usepackage[unicode,hidelinks,bookmarks=false]{hyperref}
\newcommand{\ie}{i.e.\ }
\newcommand{\cf}{cf.\ }
\newcommand{\resp}{respectively\ }
\newcommand{\Subsection}{\subsection}
\providecommand{\nobreakdash}{\nobreak\hbox{-}}
\setlength{\emergencystretch}{2em}
\multlinegap0pt
\usepackage{amssymb}
\usepackage{mathrsfs}

\DeclareMathOperator{\SL}{\mathsf{SL}}
\DeclareMathOperator{\PSL}{\mathsf{PSL}}
\DeclareMathOperator{\Gr}{\mathcal{G}}
\DeclareMathOperator{\GL}{\mathsf{GL}}
\DeclareMathOperator{\SO}{\mathsf{SO}}
\DeclareMathOperator{\Sp}{\mathsf{Sp}}
\DeclareMathOperator{\idd}{id}
\DeclareMathOperator{\tr}{tr}
\DeclareMathOperator{\spann}{span}
\DeclareMathOperator{\dist}{dist}
\DeclareMathOperator{\rSL}{SL}
\DeclareMathOperator{\Ad}{Ad}
\DeclareMathOperator{\Aut}{Aut}


\newcommand{\clA}{\mathcal{A}}
\newcommand{\F}{\mathcal{F}}
\newcommand{\N}{\mathbb{N}}
\newcommand{\R}{\mathbb{R}}
\newcommand{\RR}{\R}
\newcommand{\ZZ}{\mathbb{Z}}
\newcommand{\C}{\mathbb{C}}
\newcommand{\HH}{\mathbb{H}}
\newcommand{\sG}{\mathsf{G}}
\newcommand{\scU}{\mathscr{U}}

\newcommand{\Es}{E^s}
\newcommand{\Ec}{E^c}
\newcommand{\Eu}{E^u}
\newcommand{\Ecs}{E^{cs}}
\newcommand{\Ecu}{E^{cu}}
\newcommand{\eps}{\varepsilon}
\newcommand{\RP}{{\RR\mathbb{P}}}
\renewcommand{\a}{\mathfrak{a}}
\newcommand{\ap}{\mathfrak{a}^+}
\newcommand{\liesl}{\mathfrak{sl}_2(\RR)}

\renewcommand*{\to}{\mathchoice{\longrightarrow}{\rightarrow}{\rightarrow}{\rightarrow}}

\let\originalleft\left 
\let\originalright\right
\renewcommand{\left}{\mathopen{}\mathclose\bgroup\originalleft}
\renewcommand{\right}{\aftergroup\egroup\originalright}
%%%%%%%%%%%%
\let\oldmapsto\mapsto
\renewcommand*{\mapsto}{\mathchoice{\longmapsto}{\oldmapsto}{\oldmapsto}{\oldmapsto}}

\let\oldin\in
\renewcommand{\in}{\mathchoice{\oldin}{\oldin}{\,\oldin\,}{\,\oldin\,}}

\newcommand*{\mk}{\mkern -1mu}
\newcommand*{\Mk}{\mkern -2mu}
\newcommand*{\mK}{\mkern 1mu}
\newcommand*{\MK}{\mkern 2mu}

\hypersetup{urlcolor=purple, linkcolor=blue, citecolor=red}

\newcommand*{\relabel}{\renewcommand{\labelenumi}{(\theenumi)}}
\newcommand*{\romanenumi}{\renewcommand*{\theenumi}{\roman{enumi}}\relabel}
\newcommand*{\Romanenumi}{\renewcommand*{\theenumi}{\Roman{enumi}}\relabel}
\newcommand*{\alphenumi}{\renewcommand*{\theenumi}{\alph{enumi}}\relabel}
\newcommand*{\Alphenumi}{\renewcommand*{\theenumi}{\Alph{enumi}}\relabel}
\let\oldtilde\tilde
\renewcommand*{\tilde}[1]{\mathchoice{\widetilde{#1}}{\widetilde{#1}}{\oldtilde{#1}}{\oldtilde{#1}}}
\let\oldhat\hat
\renewcommand*{\hat}[1]{\mathchoice{\widehat{#1}}{\widehat{#1}}{\oldhat{#1}}{\oldhat{#1}}}

\numberwithin{equation}{section}
\newtheorem{theo}{Theorem}[section]
\newtheorem{prop}[theo]{Proposition}
\newtheorem{coro}[theo]{Corollary}
\newtheorem{lemm}[theo]{Lemma}
\theoremstyle{definition}\newtheorem{defi}[theo]{Definition}
\newtheorem{rema}[theo]{Remark}
\title[Robustly faithful representations into $\SL_2(\R)$ are Anosov]{Quasi-isometric free group representations into $\SL_3(\R)$: original Proposition 2.2}
\author{Le\'on Carvajales and Pablo Lessa and Rafael Potrie}
\date{Source: Annales Henri Lebesgue 8 (2025), 965--994; DOI 10.5802/ahl.252}
\begin{document}\maketitle
\noindent\textit{Editorial scope.} The counted unit is original Proposition 2.2 with its complete authored proof. The original definitions of robustness, quasi-isometric and Anosov representations and the rank-one coincidence sentence are retained as uncounted context, together with the complete original Lemma 2.1 (commutator submersion) used in the proof. Later products, examples and Barbot-representation sections are omitted. Out-of-excerpt theorem and subsection numbers are carried as hyperlinks to the publisher version; the two published bibliography entries actually cited by the retained proof are transcribed exactly.
\section{Original definitions: robust faithfulness, quasi-isometric and Anosov representations}
We will now recall the necessary definitions. We say a property of a representation $\rho:\Gamma \to \sG$ is \emph{robust} if it holds in an open neighborhood of $\rho$, in the topology of pointwise convergence. This topology is equivalent to the one induced by evaluating a representation on a given finite generating set.

Quasi-isometric representations are defined as follows. Fixing $F \subset \Gamma$ a finite symmetric generating set of $\Gamma$\footnote{Throughout, we assume that the rank of $\Gamma$ as a free group is $\geq 2$.}, we recall that the \emph{word length} of an element $\gamma \in \Gamma$ is defined by
\[
|\gamma| := \min\{ n \ge 1: \gamma \in F^n\},
\]
with the convention that $|1|$ (the word length of the neutral element) is $0$, and $F^n := \{ f_1\cdots f_n: f_1,\ldots,f_n \in F\}$. Fixing any left-invariant Riemannian metric on $\sG$ and letting $\dist$ denote the associated distance, a representation $\rho:\Gamma \to \sG$ is said to be \emph{quasi-isometric} if it is a quasi-isometric embedding of $\Gamma$, that is, there exist constants $a \ge 1$ and $b \ge 0$ such that
\begin{equation}\label{eq: QI}
a^{-1}\left|\gamma^{-1}\eta\right| - b\le \dist(\rho(\gamma),\rho(\eta)) \le a\left|\gamma^{-1}\eta\right| + b,
\end{equation}
for all $\gamma,\eta \in \Gamma$.

Finally, we recall the definition of Anosov representations. We fix a Cartan subspace $\a$ of $\sG$, a Weyl chamber $\ap\subset\a$, and a corresponding Cartan projection $a:\sG \to \ap$. A representation is said to be \emph{Anosov} if there exist constants $c,d > 0$ and a simple root $\alpha$ such that
\begin{equation}\label{eqanosov}
c |\gamma| - d \le \alpha(a(\rho(\gamma))),
\end{equation}
for all $\gamma \in \Gamma$. In that case we say that $\rho$ is $\{\alpha\}$-\emph{Anosov}.

From Equation~\eqref{eqanosov} one readily sees that Anosov representations are quasi-isometric. While in rank one these two notions coincide, in higher rank being quasi-isometric is not a robust property (see Subsection~\href{https://ahl.centre-mersenne.org/item/AHL_2025__8_965_0/}{1.3.1} below).

\section{Products of \texorpdfstring{$\SL_2$}{SL\_2}'s}\label{sec: products}

\subsection{Robust faithfulness in $\SL_2(\R)$}\label{ss.SL2}

Here we show that robustly faithful representations of a free group into $\SL_2(\R)$ are Anosov (note that this does not follow from Sullivan's Theorem~\href{https://ahl.centre-mersenne.org/item/AHL_2025__8_965_0/}{1.3}). This will be used in the proof of Theorem~\href{https://ahl.centre-mersenne.org/item/AHL_2025__8_965_0/}{1.4}.

\begin{lemm}\label{lem: commutator is submersion}
Let $g_0$ and $h_0$ be non commuting hyperbolic elements in $\SL_2(\R)$.
Then the commutator
\[
[\cdot,\cdot]:\SL_2(\R)\times\SL_2(\R)\to\SL_2(\R); [g,h]:=ghg^{-1}h^{-1}
\]
is a submersion at $(g_0,h_0)$.
\end{lemm}


\begin{proof}
This is a direct computation, we include a proof for completeness.


Let $\liesl$ be the Lie algebra of $\SL_2(\R)$, that is, the vector space of traceless $2\times 2$ real matrices. For $g\in\SL_2(\R)$, the tangent space $T_g\SL_2(\R)$ is $g\cdot\liesl$.

Let $X,Y\in\liesl$.
The derivative of $[\cdot,\cdot]$ at $(g_0,h_0)$ evaluated on the pair $(g_0\cdot X,h_0\cdot Y)$ is given by
\[
[g_0,h_0]\cdot \left(\Ad_{h_0g_0h_0^{-1}}(X)+\Ad_{h_0g_0}(Y)-\Ad_{h_0g_0}(X)-\Ad_{h_0}(Y)\right),
\]
where $\Ad:\SL_2(\R)\to\Aut(\liesl)$ is the adjoint representation.
In particular, it suffices to show that the linear map
\[
(X,Y)\mapsto \Ad_{g_0}\left(\Ad_{h_0^{-1}}(X)-X\right)+\left(\Ad_{g_0}(Y)-Y\right)
\]
is surjective.

Without loss of generality we may suppose that $g_0$ is diagonal, with eigenvalues $\mu\neq \pm 1$ and $\mu^{-1}$. Then $\Ad_{g_0}$ is diagonalizable on the basis $\left\{\left(
\begin{smallmatrix}
1 & 0 \\
0 & -1
\end{smallmatrix}\right),\left(
\begin{smallmatrix}
0 & 1 \\
0 & 0
\end{smallmatrix}\right),\left(
\begin{smallmatrix}
0 & 0 \\
1 & 0
\end{smallmatrix}\right)\right\}$, with respective eigenvalues $1$, $\mu^2$ and $\mu^{-2}$.
The image of $Y\mapsto \Ad_{g_0}(Y)-Y$ is therefore
\[
W_{g_0}:=\spann\left\{\left(
\begin{matrix}
0 & 1 \\
0 & 0
\end{matrix}\right),\left(
\begin{matrix}
0 & 0 \\
1 & 0
\end{matrix}\right)\right\}.
\]
Moreover, as $h_0$ is hyperbolic but not diagonal (because it does not commute with $g_0$), it follows easily that $W_{g_0}$ is different from the image of
\[
X\mapsto \Ad_{g_0}\left(\Ad_{h_0^{-1}}(X)-X\right),
\]
which is again two-dimensional.
This completes the proof of Lemma~\ref{lem: commutator is submersion}.
\end{proof}

\begin{prop}\label{prop:robustfaithfulSL2}
Let $\Gamma$ be a finitely generated non-abelian free group and $\rho: \Gamma \to \SL_2(\R)$ a robustly faithful representation. Then $\rho$ is Anosov (equivalently, quasi-isometric).
\end{prop}

\begin{proof}
We first observe that $\rho$ must be discrete.
Otherwise, the Euclidean closure $\overline{\rho(\Gamma)}$ of $\rho(\Gamma)$ would be either $\SL_2(\R)$ or contained in a (virtually) solvable subgroup of $\SL_2(\R)$. The latter case is ruled out, as it would imply that $\rho$ has a non trivial kernel. Hence $\overline{\rho(\Gamma)} = \SL_2(\R)$ and there is an element $\gamma\in\Gamma\setminus\{1\}$ so that the trace $\tr(\rho(\gamma))$ belongs to the interval $(-2,2)$. Since the trace function $\rho' \mapsto \tr(\rho'(\gamma))$ is a non-constant polynomial in the entries of the image of a free generating set of $\Gamma$, it must take a value of the form 
$2 \cos(\frac{k \pi}{n})$ for some integers $k$ and $n$ and arbitrarily small perturbations $\rho'$ of $\rho$. Thus we would get a small perturbation of $\rho$ which is not faithful contradicting our assumption.

Now that we know that $\rho(\Gamma)$ is discrete, observe that the induced representation $\Gamma\to\PSL_2(\R)$ is also discrete. It is moreover robustly faithful and, to finish the proof, it suffices to show that this induced representation is quasi-isometric. By abuse of notations, we still denote this induced representation by $\rho$.

Consider the hyperbolic surface $\HH^2/\rho(\Gamma)$, which is geometrically finite as $\Gamma$ is finitely generated (see e.g.~\cite[Theorem~4.6.1]{skatok}). We will show that every element in $\rho(\Gamma)$ is hyperbolic, and this will finish the proof (see e.g.~\cite[Theorems~4.8 \& 4.13 and Corollary~4.17]{dalbogeodesichorocyclic}).



Suppose by contradiction that there is some parabolic element corresponding to some cusp in $\HH^2/\rho(\Gamma)$. There are two cases, namely, either $\HH^2/\rho(\Gamma)$ has genus $g=0$, or $g>0$.

If $g=0$, we may suppose that the cusp corresponds to a free generator of $\Gamma$.
It can then be perturbed to a finite order element in $\PSL_2(\R)$, contradicting the fact that $\rho$ is robustly faithful. On the other hand, if $g>0$ then the cusp corresponds to an element $\gamma$ of $\Gamma$ so that there is a free generating set
\[
\{a_1,b_1,\dots,a_g,b_g,c_1,\dots,c_n\}\subset\Gamma
\]
such that
\[
\gamma=[a_1,b_1]\dots[a_g,b_g]c_1\dots c_n.
\]
We can assume that both $\rho(a_1)$ and $\rho(b_1)$ are hyperbolic, otherwise we conclude as in the previous case. Since $\rho$ is faithful, by Lemma~\ref{lem: commutator is submersion} we may perturb $\rho$ only perturbing the image of $a_1$ and $b_1$, in such a way that $\rho(\gamma)$ becomes a finite order element of $\PSL_2(\R)$. This finishes the proof of Proposition~\ref{prop:robustfaithfulSL2}.
\end{proof}

\begin{thebibliography}{99}
\bibitem[Kat92]{skatok} Svetlana Katok, \emph{Fuchsian groups}, University of Chicago Press, 1992.
\bibitem[Dal11]{dalbogeodesichorocyclic} Fran\c{c}oise Dal'Bo, \emph{Geodesic and horocyclic trajectories}, Universitext, EDP Sciences; Springer, 2011, translation from the french by the author.
\end{thebibliography}
\end{document}
